% Enumeration of (translation, shift) pairs in GbShifts:
% the Flat and Full searches.
%
% Self-contained: compiles with pdflatex, needs only amsmath/amssymb.
% To lift into a larger document, take everything between \begin{document}
% and \end{document} and drop the \section* levels one step.

\documentclass[11pt]{article}
\usepackage{amsmath}
\usepackage{amssymb}
\usepackage[margin=1in]{geometry}

\newcommand{\Lat}[1]{\mathcal{#1}}
\newcommand{\A}{\Lat{A}}
\newcommand{\B}{\Lat{B}}
\newcommand{\C}{\Lat{C}}
\newcommand{\D}{\Lat{D}}
\newcommand{\LA}{\Lambda_{\A}}
\newcommand{\LB}{\Lambda_{\B}}
\newcommand{\nhat}{\hat{\mathbf{n}}}
\newcommand{\tmax}{t_{\max}}
\newcommand{\tperp}{t_{\perp\max}}
\newcommand{\sperp}{s_{\perp\max}}

\title{Enumeration of translation--shift pairs for grain-boundary mesostates}
\date{}

\begin{document}
\maketitle

\section*{Notation}

Let $\A$ and $\B$ be the two grain lattices, $\C$ their coincidence-site
lattice (CSL) and $\D$ the displacement-shift-complete lattice (DSCL), so that
$\C \subset \A,\B \subset \D$. Write

\begin{itemize}
  \item $b = \lVert \mathbf{a}_1 \rVert$, the shortest lattice vector of $\A$,
        used as the unit for every length tolerance;
  \item $\mathsf{D}$ for the DSCL basis matrix, so every translation is
        $\mathbf{t} = \mathsf{D}\mathbf{n}$ with $\mathbf{n} \in \mathbb{Z}^{3}$;
  \item $\LA, \LB$ for the shift tensors, which satisfy
        \begin{equation}
          \LA + \LB = \mathsf{I},
          \qquad
          \LA \mathbf{t} \in \B,
          \qquad
          \LB \mathbf{t} \in \A
          \qquad \text{for all } \mathbf{t} \in \D ;
          \label{eq:shift-tensors}
        \end{equation}
  \item $\mathbf{n}_{\C}$ for the boundary normal expressed as a reciprocal
        lattice direction of $\C$, $\nhat$ for its unit Cartesian form, and
        $d(\mathbf{n}_{\C})$ for the corresponding interplanar spacing;
  \item $\mathbf{p}_1, \mathbf{p}_2$ for the two in-plane CSL vectors that span
        the boundary, and $\mathbf{c}_0$ for the out-of-plane member of the
        plane-parallel CSL basis about $\mathbf{n}_{\C}$.
\end{itemize}

For a basis $\mathsf{V} = [\,\mathbf{v}_1 \mid \mathbf{v}_2 \mid \mathbf{v}_3\,]$
and an offset $\boldsymbol{\sigma} \in \mathbb{R}^{3}$, define the reduction
$\bmod(\,\cdot\,;\mathsf{V},\boldsymbol{\sigma})$ as the representative whose
coordinates in $\mathsf{V}$ lie in $[\sigma_i, \sigma_i + 1)$:
\begin{equation}
  \bmod(\mathbf{x};\mathsf{V},\boldsymbol{\sigma})
  = \mathbf{x} - \mathsf{V}
    \left\lfloor \mathsf{V}^{-1}\mathbf{x} - \boldsymbol{\sigma} \right\rfloor ,
  \label{eq:modulo}
\end{equation}
the floor acting componentwise.

\section*{Construction common to both searches}

\subsection*{The mesostate CSL cell}

Both searches reduce shifts into the same cell: the two in-plane CSL vectors,
together with however many CSL planes are needed to accommodate the requested
layer thickness,
\begin{equation}
  m = \max\!\left(1,\;
      \left\lfloor \frac{\sperp\, b}{d(\mathbf{n}_{\C})} \right\rfloor \right),
  \qquad
  \mathsf{C} = [\, m\,\mathbf{c}_0 \mid \mathbf{p}_1 \mid \mathbf{p}_2 \,].
  \label{eq:cell}
\end{equation}
Let $P = \C \cap \mathsf{C}$ denote the CSL sites inside that cell. The
reduction is taken modulo this \emph{mesostate} cell rather than the primitive
CSL cell: reducing modulo the primitive cell confines every shift to a fraction
of the box and leaves most of the boundary without nodes.

\subsection*{The shift}

For a translation $\mathbf{t}$ and a CSL site $\mathbf{p} \in P$,
\begin{equation}
  \mathbf{s} = \LA \mathbf{t} + \mathbf{p} - \tfrac{1}{2}\mathbf{t},
  \label{eq:shift}
\end{equation}
followed by reduction into the cell and the acceptance test
\begin{equation}
  \left\lvert \mathbf{s} \cdot \nhat \right\rvert \le \tfrac{1}{2}\,\sperp\, b .
  \label{eq:accept}
\end{equation}
The reduction uses the offset $\boldsymbol{\sigma} = (-\tfrac12, 0^{-}, 0^{-})$,
so the out-of-plane coordinate is centred on zero and
$\mathbf{s}\cdot\nhat \in [-\tfrac12 m\,d, \tfrac12 m\,d)$, while the two
in-plane coordinates run over $[0,1)$.

Pairing each translation with \emph{every} $\mathbf{p} \in P$ is essential. Were
$\mathbf{s}$ a function of $\mathbf{t}$ alone, all sites generated by one
translation would fold onto a single lattice point and every combination of them
would be rejected as a clash; in particular a flat boundary, on which many sites
share one translation, would be unreachable.

\subsection*{Why every DSCL translation is admissible}

The two nodes contributed by a pair are $\mathbf{x}_{\A} = \mathbf{s} -
\tfrac12\mathbf{t}$ and $\mathbf{x}_{\B} = \mathbf{s} + \tfrac12\mathbf{t}$.
Substituting~\eqref{eq:shift} and using~\eqref{eq:shift-tensors},
\begin{align}
  \mathbf{x}_{\A}
    &= \LA\mathbf{t} + \mathbf{p} - \mathbf{t}
     = -\LB\mathbf{t} + \mathbf{p} \;\in\; \A ,
  \label{eq:xA}\\[2pt]
  \mathbf{x}_{\B}
    &= \LA\mathbf{t} + \mathbf{p} \;\in\; \B .
  \label{eq:xB}
\end{align}
These hold for \emph{any} $\mathbf{t} \in \D$ and any $\mathbf{p} \in \C$, with
no restriction on the direction of $\mathbf{t}$, and reduction modulo
$\mathsf{C}$ preserves them because the cell vectors are CSL vectors and hence
belong to both lattices. Widening the set of translations therefore cannot
violate the lattice-membership requirement downstream.

\section*{The flat search}

The translations are drawn from a parallelepiped whose edges are DSCL vectors
aligned with the CSL box directions. Starting from the cell
edges $\mathbf{v}_i$ of~\eqref{eq:cell}, reoriented toward orthogonality, each is
pulled into the DSCL and scaled,
\begin{equation}
  \mathbf{d}_i = \operatorname{latticeDirection}_{\D}(\mathbf{v}_i),
  \qquad
  f_i = \max\!\left(1,\;
        \left\lfloor \frac{\tmax\, b}{\lVert \mathbf{d}_i \rVert} \right\rfloor
        \right),
  \label{eq:flat-edges}
\end{equation}
giving the translation domain
\begin{equation}
  T_{\text{flat}}
  = \left\{\, \mathbf{t} \in \D \cap \Pi \;:\;
      \lVert \mathbf{t} \rVert \le \tmax\, b \,\right\},
  \qquad
  \Pi = \left\{ \textstyle\sum_i \xi_i f_i \mathbf{d}_i
        \;:\; \xi_i \in [-\tfrac12, \tfrac12) \right\}.
  \label{eq:flat-domain}
\end{equation}
Each $\mathbf{t}$ is reduced into $\Pi$ before the norm test, then
combined with every $\mathbf{p} \in P$ through~\eqref{eq:shift}, reduced modulo
$\mathsf{C}$, and accepted by~\eqref{eq:accept}.

Because the domain is built from the CSL box structure rather than from
geometry, the surviving translations are those commensurate with that structure,
and the resulting CSL shifts lie in the boundary plane. Only flat mesostates are
reachable.

\section*{The full search}

The parallelepiped is abandoned. Translations are enumerated directly in integer
coordinates over an exact bounding box for the ball: from
$\mathbf{t} = \mathsf{D}\mathbf{n}$,
\begin{equation}
  n_i = \left(\mathsf{D}^{-1}\right)_{i\cdot}\!\cdot \mathbf{t}
  \quad\Longrightarrow\quad
  \lvert n_i \rvert
  \le \left\lVert \left(\mathsf{D}^{-1}\right)_{i\cdot} \right\rVert
      \lVert \mathbf{t} \rVert
  \le \left\lVert \left(\mathsf{D}^{-1}\right)_{i\cdot} \right\rVert \tmax\, b ,
  \label{eq:bound}
\end{equation}
so it suffices to scan $n_i \in [-N_i, N_i]$ with
\begin{equation}
  N_i = \left\lceil
        \left\lVert \left(\mathsf{D}^{-1}\right)_{i\cdot} \right\rVert \tmax\, b
        \right\rceil + 1 .
  \label{eq:scanlimit}
\end{equation}
Two geometric filters then define the domain,
\begin{equation}
  T_{\text{full}}
  = \left\{\, \mathbf{t} \in \D \;:\;
      \lVert \mathbf{t} \rVert \le \tmax\, b
      \;\text{ and }\;
      \lvert \mathbf{t} \cdot \nhat \rvert \le \tperp\, b \,\right\},
  \label{eq:full-domain}
\end{equation}
the slab being inert whenever $\tperp \ge \tmax$, since
$\lVert\mathbf{t}\rVert \le \tmax b$ already implies
$\lvert \mathbf{t}\cdot\nhat \rvert \le \tmax b$.

\subsection*{Two-stage reduction}

The shift is formed by~\eqref{eq:shift} exactly as before, but the reduction has
a second stage. Write $\mathsf{P} = [\,\mathbf{p}_1 \mid \mathbf{p}_2\,]$. After
the reduction modulo $\mathsf{C}$, solve the least-squares problem
\begin{equation}
  \boldsymbol{\alpha}
  = \arg\min_{\boldsymbol{\alpha} \in \mathbb{R}^{2}}
    \left\lVert \mathsf{P}\boldsymbol{\alpha} - \mathbf{s} \right\rVert ,
  \qquad
  \mathbf{s} \leftarrow
  \mathbf{s} - \mathsf{P}
  \left\lfloor \boldsymbol{\alpha} + \epsilon \right\rfloor ,
  \label{eq:inplane-fold}
\end{equation}
before applying the acceptance test~\eqref{eq:accept}.

This stage is not redundant. The vector $\mathbf{c}_0$ returned by the
plane-parallel basis is sheared: it carries in-plane components, so a shift
reduced in the basis $\mathsf{C}$ can still lie up to a quarter of a period
outside the requested box. Subtracting whole multiples of $\mathbf{p}_1$ and
$\mathbf{p}_2$ moves the shift onto an equivalent site, and because
\begin{equation}
  \mathbf{p}_1 \cdot \nhat = \mathbf{p}_2 \cdot \nhat = 0 ,
  \label{eq:inplane-orthogonal}
\end{equation}
the quantity $\mathbf{s} \cdot \nhat$ is invariant under~\eqref{eq:inplane-fold}.
The fold therefore cannot disturb the acceptance test applied around it. The
offset $\epsilon$ before the floor guards the boundary case: a shift belonging at
coordinate $1$, the periodic image of $0$, routinely lands a rounding error
below it and would otherwise be recorded as a distinct site on the far face of
the box.

The surviving pairs are ordered by $\lVert \mathbf{t} \rVert$, with ties broken
lexicographically on integer coordinates, so that the pair
$(\mathbf{0},\mathbf{0})$ heads the list and signature indices are reproducible.
Optionally, only the shortest translation among those sharing a shift is kept.

\section*{How the two searches differ}

The shift machinery is common to both: equation~\eqref{eq:shift} for the shift,
the cell~\eqref{eq:cell} it is reduced into, and the
criterion~\eqref{eq:accept} for accepting it are identical. The full search
departs from the flat one in exactly two places.

First, the translation domain. The flat search
takes~\eqref{eq:flat-domain}, a lattice-aligned parallelepiped trimmed by a
norm test, so its translations inherit the CSL box structure and their CSL
shifts lie in the boundary plane. The full search
takes~\eqref{eq:full-domain}, an unbiased geometric neighbourhood --- a ball
intersected with a slab --- enumerated through the exact integer
bound~\eqref{eq:scanlimit}, and introduces the additional parameter $\tperp$.
Dropping the alignment is what admits translations whose CSL shift retains an
out-of-plane component, $\mathbf{s} \cdot \nhat \neq 0$, and hence what makes
non-flat mesostates reachable.

Second, the reduction. The flat search reduces once, modulo
$\mathsf{C}$. The full search follows that with the in-plane
fold~\eqref{eq:inplane-fold}, which is what keeps the newly admitted shifts on
distinct, correctly-boxed sites instead of scattering them outside the cell,
while leaving $\mathbf{s}\cdot\nhat$ --- and therefore the acceptance
test --- untouched.

Two consequences follow. The pair count is governed by $\tmax$ and $\tperp$
through~\eqref{eq:full-domain}, while $\sperp$ alone decides how far off the
boundary plane a shift may sit, through~\eqref{eq:cell}
and~\eqref{eq:accept}. Setting $\sperp \to 0$ collapses the admissible layer to
the boundary plane, and the full search degenerates to a flat-only enumeration.

\end{document}
